% Các tệp trong appendix/src/ được build.sh sao chép từ kho mã nguồn (projects/fgtn_ava) trước khi biên dịch.
% Phần thân không tham chiếu tới nhãn trong tệp này, nên tắt phụ lục (\inphulucfalse) vẫn biên dịch được.

\chapter{Cấu hình huấn luyện}
\label{app:config}

Phụ lục này in nguyên văn hai tệp cấu hình MMAction2 trong kho mã \path{projects/fgtn_ava}. Các cấu hình ablation A0--A7 được tạo từ hai tệp này bằng một khóa ghi đè duy nhất cho mỗi cấu hình (\cref{app:ablation-map}).

\section{Cấu hình chỉ dùng tầng 4 --- cơ sở của A0--A6}
\label{app:config-single}

Cấu hình này không dùng tầng 3 và ACRN. Tệp \path{ava22_swin-tiny_fgtn_32x2x1-10e.py}: Video Swin-Tiny tiền huấn luyện Kinetics-400, khối Hybrid FGTN-Mamba trên tầng 4, bộ trích đặc trưng vùng \code{SingleRoIExtractor3D} với \code{featmap\_stride = 32} và head sigmoid đa nhãn.
\lstinputlisting[language=Python]{appendix/src/ava22_swin-tiny_fgtn_32x2x1-10e.py}

\section{Cấu hình pipeline đầy đủ --- A8}
\label{app:config-full}

Tệp \path{ava22_swin-tiny_fgtn-ms-acrn-focal_32x2x1-10e.py} kế thừa cấu hình chỉ dùng tầng 4 và bổ sung ba thành phần: đặc trưng đa tỉ lệ (tầng 3 và tầng 4, \code{in\_channels = (384, 768)}), ACRN (\code{ACRNHead} làm \code{shared\_head}) và head focal có trọng số lớp.
\lstinputlisting[language=Python]{appendix/src/ava22_swin-tiny_fgtn-ms-acrn-focal_32x2x1-10e.py}

\chapter{Mã nguồn khối đề xuất}
\label{app:code}

\section{Hybrid FGTN-Mamba neck}

Tệp \path{fgtn_neck.py}: khối hai nhánh (Bi-SSM và W-MHSA) hợp nhất bằng cổng sigmoid, ngăn xếp riêng cho từng tầng đặc trưng. Tham số \code{branch\_mode} chọn một trong bốn chế độ \code{hybrid}, \code{mamba\_only}, \code{attn\_only}, \code{no\_gate} dùng cho ablation.
\lstinputlisting[language=Python]{appendix/src/fgtn_neck.py}

\section{Head focal + trọng số lớp}

Tệp \path{heads.py}: head sigmoid đa nhãn với focal loss ($\gamma = 2$) nhân trọng số lớp.
\lstinputlisting[language=Python]{appendix/src/heads.py}

\section{Tính trọng số lớp theo số mẫu hiệu dụng}

Tệp \path{make_class_weight.py}: tính $w_c \propto (1-\beta)/(1-\beta^{n_c})$ với $\beta = 0{,}999$, chuẩn hóa trung bình 1 và cắt trên tại 10, ghi ra tệp JSON cho khóa \code{class\_weight} của head.
\lstinputlisting[language=Python]{appendix/src/make_class_weight.py}

\section{Kiểm tra trên CPU}

Tệp \path{selfcheck.py}: kiểm tra kích thước tensor của neck ở bốn chế độ của tham số \code{branch\_mode}, của head và của hàm tính trọng số lớp mà không cần GPU hay MMAction2.
\lstinputlisting[language=Python]{appendix/src/selfcheck.py}

\chapter{Ánh xạ cấu hình ablation sang mã nguồn}
\label{app:ablation-map}

\Cref{tab:ablation-map} cho biết mỗi cấu hình A0--A8 của \cref{tab:ablation} được tạo từ tệp cấu hình nào và khóa ghi đè nào. Mỗi cặp so sánh chỉ khác nhau ở đúng một khóa.

\begin{longtable}{@{}>{\raggedright\arraybackslash}p{1.0cm}>{\raggedright\arraybackslash}p{4.6cm}>{\raggedright\arraybackslash}p{8.6cm}@{}}
\caption{Ánh xạ cấu hình ablation A0--A8 sang tệp cấu hình và khóa ghi đè}
\label{tab:ablation-map}\\
\toprule
\textbf{Mã} & \textbf{Cấu hình} & \textbf{Tệp cấu hình và khóa ghi đè (\code{--cfg-options})}\\
\midrule
\endfirsthead
\toprule
\textbf{Mã} & \textbf{Cấu hình} & \textbf{Tệp cấu hình và khóa ghi đè (\code{--cfg-options})}\\
\midrule
\endhead
A0 & Tầng 4, không neck & Chỉ dùng tầng 4; \path{model.neck=None}\\
A1 & + temporal MHSA toàn cục & Chỉ dùng tầng 4; \path{model.neck.branch_mode=attn_only}, \path{model.neck.window_size=31} (cửa sổ phủ toàn bộ $T = 16$)\\
A2 & + SSM một chiều & Chỉ dùng tầng 4; \path{branch_mode=mamba_only} và đặt \path{bimamba_type=none} trong \path{fgtn_neck.py} (chưa có khóa cấu hình)\\
A3 & + Bi-SSM & Chỉ dùng tầng 4; \path{model.neck.branch_mode=mamba_only}\\
A4 & + W-MHSA cục bộ ($W = 5$) & Chỉ dùng tầng 4; \path{model.neck.branch_mode=attn_only}\\
A5 & Hai nhánh, cộng trực tiếp 0,5 / 0,5 & Chỉ dùng tầng 4; \path{model.neck.branch_mode=no_gate}\\
A6 & Hybrid có cổng sigmoid (tầng 4) & Chỉ dùng tầng 4, giữ mặc định \path{branch_mode=hybrid}\\
A7 & A6 trên tầng 3 + 4 & Đầy đủ; \path{model.roi_head.shared_head=None} (bỏ ACRN)\\
A8 & A7 + ACRN (pipeline đầy đủ) & Đầy đủ, giữ nguyên\\
\bottomrule
\end{longtable}

Trong cặp A6 $\rightarrow$ A7, hàm mất mát cần giữ giống nhau; khi so với A6 dùng BCE, đặt thêm \path{model.roi_head.bbox_head.focal_gamma=0.0} cho A7. Trọng số lớp được đặt trên khóa của head (\path{model.roi_head.bbox_head.class_weight}); khóa \path{class_weight} ở cấp cao nhất đã được đọc khi nạp cấu hình nên ghi đè nó không có tác dụng.

\chapter{Bảng truy nguồn số liệu}
\label{app:traceability}

Mỗi con số thực nghiệm trong luận văn được ghi kèm nguồn trong kho mã, trong nhật ký huấn luyện và đánh giá của luận văn hoặc trong tài liệu được trích dẫn (\cref{tab:trace}). Các đường dẫn \path{logs/} và \path{work_dirs/} là thư mục nhật ký và kết quả của các lượt chạy, được lưu kèm kho mã; mã nguồn nằm trong \path{projects/fgtn_ava/}.

\begin{longtable}{@{}>{\raggedright\arraybackslash}p{4.1cm}>{\raggedright\arraybackslash}p{2.7cm}>{\raggedright\arraybackslash}p{7.4cm}@{}}
\caption{Truy nguồn các số liệu trong luận văn}
\label{tab:trace}\\
\toprule
\textbf{Số liệu} & \textbf{Giá trị} & \textbf{Nguồn}\\
\midrule
\endfirsthead
\toprule
\textbf{Số liệu} & \textbf{Giá trị} & \textbf{Nguồn}\\
\midrule
\endhead
Phân bố lượt gán nhãn của bộ dữ liệu phân xưởng & \cref{tab:factory-stats} (24.614 lượt) & Thống kê tệp nhãn của bộ dữ liệu phân xưởng\\
Phép chia theo video nguồn; rò rỉ của phân chia gốc & \cref{tab:vsplit}; 2.891/3.831 khung hình khóa & \path{make_hpd_split.py} (seed 42) chạy trên tệp nhãn gốc; đối chiếu cặp (video, khung hình khóa) giữa \path{train.csv} và \path{val.csv}\\
Cấu hình huấn luyện trên dữ liệu phân xưởng & \cref{tab:factory-config} & \path{hpd_swin-tiny_fgtn_debias_wandb.py}, \path{hpd_swin-tiny_fgtn-ms-acrn_debias_wandb.py}; dòng lệnh và learning rate trong nhật ký huấn luyện\\
AP theo lớp, mAP theo epoch của A0, A6, A8, A6-gốc, A6+AVA & \cref{tab:main-results,tab:epochs,tab:ablation-results} & Nhật ký console \path{logs/debias.log}, \path{logs/a8.log}, \path{logs/train_hpd.log}, \path{logs/night.log} (khối ``Per-class results'' sau mỗi epoch kiểm định) và nhật ký MMEngine trong \path{work_dirs/hpd*/<cấu hình>/<thời điểm>/}\\
Thăm dò che người, che bối cảnh, bỏ chuyển động & \cref{tab:probe} & \path{probe_scene.py}; dòng \code{PROBE\_SUMMARY} trong \path{logs/night.log} và \path{logs/probe_A5_ep2.log}\\
Chỉ số đầu--cuối với YOLO26-x & \cref{tab:operating} & \path{demo_video.py --dump} trên 23 video kiểm định (tệp JSON trong \path{work_dirs/torch_val/<mô hình>/}); tính bằng \path{eval_e2e.py} (ghép với nhãn bằng \code{score\_frame} của \path{score_val_demos.py}); \cref{fig:demo} vẽ bằng \path{render_demo_fig.py}\\
Đánh giá logic cảnh báo & \cref{tab:alert-eval} & \path{alert.py} chạy trên cùng các tệp JSON, $m/n = 2/3$, cooldown 5~s, ngưỡng 0,5\\
Thời gian suy luận, bộ nhớ GPU, thời gian mỗi epoch & \cref{tab:cost} & Trung vị trường \code{time} và giá trị lớn nhất của \code{memory} ở các dòng \code{Epoch(val)}/\code{Epoch(train)} trong nhật ký MMEngine; khoảng cách giữa các dòng kết quả kiểm định\\
mAP trên tập con AVA~v2.2 (A0, A6) & \cref{tab:ava-subset} & Nhật ký \path{logs/train_ablation.log} (\path{work_dirs/subset/A0_noneck}, \path{work_dirs/subset/A5_hybrid}); AP theo lớp trong cùng nhật ký\\
Bước giảm mẫu tầng 4 của Swin-T & 32 & Patch $(2,4,4)$ và ba lần gộp patch; \path{fgtn_ava/README.md} ghi $(B, 768, 16, H/32, W/32)$\\
Số tham số ACRN, neck tầng 3 & $\approx 14{,}6$M; 3,57M & Tính từ kích thước lớp trong \code{ACRNHead} và \code{FGTNNeckAVA(384)}\\
Mốc A*, STAR, OV-STAD, ST-CLIP & \cref{tab:reference} & Bài báo gốc \cite{kim2024astar,gritsenko2024star,wu2024ovstad,huang2025stclip}\\
\bottomrule
\end{longtable}

\chapter{Lệnh cài đặt, huấn luyện và kiểm tra}
\label{app:repro}

Các lệnh dưới đây chạy từ thư mục gốc của kho mã (khóa \code{custom\_imports} trong cấu hình cần đường dẫn \path{projects.fgtn_ava}). Với bộ dữ liệu phân xưởng, chỉ thay đường dẫn dữ liệu, tệp nhãn và số lớp trong cấu hình; khóa ghi đè của từng cấu hình ablation giữ nguyên như \cref{tab:ablation-map}.
\begin{lstlisting}[language=bash,numbers=none]
# 1. Environment (see projects/fgtn_ava/README.md)
pip install -U openmim && mim install mmengine "mmcv>=2.0.0,<2.2.0" mmdet "mmaction2>=1.2.0"
pip install causal-conv1d
pip install git+https://github.com/OpenGVLab/video-mamba-suite.git#subdirectory=mamba

# 2. CPU self-check (no mmaction / GPU required)
python projects/fgtn_ava/selfcheck.py

# 3. Stage-4-only Hybrid FGTN-Mamba (A6);
#    A0-A5 add one override each (Table C.1)
SS=projects/fgtn_ava/ava22_swin-tiny_fgtn_32x2x1-10e.py
# A6
mim train mmaction2 $SS --gpus 8 --launcher pytorch
# A5
mim train mmaction2 $SS --gpus 8 --launcher pytorch \
    --cfg-options model.neck.branch_mode=no_gate

# 4. Class weights (effective number of samples)
python projects/fgtn_ava/make_class_weight.py \
    --ann data/ava/annotations/ava_train_v2.2.csv \
    --out data/ava/annotations/class_weight_train.json
W=data/ava/annotations/class_weight_train.json

# 5. Full pipeline: multi-scale + ACRN + focal
#    + class weights (A8); A7 = A8 without ACRN
MS=projects/fgtn_ava/ava22_swin-tiny_fgtn-ms-acrn-focal_32x2x1-10e.py
# A8
mim train mmaction2 $MS --gpus 8 --launcher pytorch \
    --cfg-options model.roi_head.bbox_head.class_weight=$W
# A7
mim train mmaction2 $MS --gpus 8 --launcher pytorch \
    --cfg-options model.roi_head.bbox_head.class_weight=$W \
    model.roi_head.shared_head=None
\end{lstlisting}
